\chapter{Conjugate angular coordinates and spectral decomposition}
\label{chap:canales}

The passage from the arithmetic closure of \(\alpha\) to the dimensional
hierarchy is effected through two angular coordinates. The first follows
directly from \(\alphaHMT\); the second is the oriented phase produced by
the regional analytic realization of the quintuple microfiber. Both reach
Dimensional Mechanics already constructed, before the exponential chart
that will generate the global moments is introduced.

\section{Transfer of the angular coordinates derived in HMT}

Parts I--III constructed distinct dependencies with a common root in
\(\alpha_{\rm HMT}\). The pre-return character follows from
\[
(\alpha_{\rm HMT},\varphi;\,9,5,4,12,54,\text{ graded rule})
\longrightarrow H_5,
\]
while the angular coordinate and the determinant follow from
\[
\alpha_{\rm HMT}\longrightarrow
A=1000\alpha_{\rm HMT}\,{}^\circ,
\qquad
(\alpha_{\rm HMT},\pi)\longrightarrow
D_A=e^{-100\pi\alpha_{\rm HMT}/9}.
\]
Through a distinct dependence, the regional microfiber fixes
\(\mathcal F_\pi\to\rho_\pi=169\), and with it
\[
\mathcal C_\pi(\alpha_{\rm HMT})
=169\alpha_{\rm HMT}^{6}
+\frac{D_A\alpha_{\rm HMT}^{7}}
       {1-D_A\alpha_{\rm HMT}}.
\]
The branch of \(H_5\) and the regional correction converge only at the
following step:
\[
\begin{aligned}
y_*&=\frac{\pi}{90}(H_5+\alpha_{\rm HMT})
     -\mathcal C_\pi(\alpha_{\rm HMT}),\\
C^*&=\frac{180}{\pi}y_*.
\end{aligned}
\]
Henceforth, \(A\) denotes the direct angular coordinate and \(C^*\) the
conjugate regional angular coordinate; these designations fix their types
without introducing a second selection rule.
Consequently, \(H_5\) is not an output of \(A\). Here \(y_*\) is the analytic
coordinate in radians and \(C^*\) its sexagesimal coordinate. Dimensional
Mechanics receives these already determined outputs; it does not redefine them.
The coordinate \(A\) is not a direct input to the formula of \(C^*\): the two
objects are coupled for the first time in \((A,C^*)\mapsto q_\pm\). The equality
\(\alpha_{\rm HMT}=A/1000\) is solely the inverse change of coordinate.
In this Part we retain the notations
\[
\begin{aligned}
 A&=7.297352569283800997285105472380663^\circ,\\
 \Cstar&=2.123738338968645724261951380393\ldots{}^\circ.
\end{aligned}
\]
The rounding used in the tables is
\(2.123738338968646^\circ\). The integer \(169\), the determinant \(D_A\), and
the regional recurrence belong to the preceding proof; neither any
SI comparison nor the hexad--octad functional intervenes in their generation.

The
angular identity
\[
 \etaRet=\frac{\Cstar}{2}-\alphaAngle,
 \qquad
 \alphaAngle:=\alphaHMT\,{\rm deg},
\]
defines the coordinate angular adimensional posterior to the return; does not  employed
for generate \(\Cstar\). Equivalently,
a time performed the construction regional,
\[
 \Cstar=2(\etaRet+\alphaAngle).
\]
This last equality is the inverse form of the same relation and not a
second selection rule for \(C^*\).

The regional realization corrects \(H_5\) before producing \(\etaRet\); therefore,
\(H_5\), \(\etaRet\), and \(\Cstar\) have distinct mathematical types. The symbol
\(\hbarHMT\) is reserved for the magnitude on the action line obtained only
after applying the scale \(10^{-34}\mathcal U_S\). Comparison with an SI mantissa belongs
to external metrological control and leaves the preceding construction invariant.

\section{Internal phases and exponential chart}

The abstract phases associated with the two angles are
\[
 a=\frac A{360},\qquad c=\frac{\Cstar}{360}
 \quad\text{in }\RR/\ZZ.
\]
Before applying a real exponential, the canonical representatives
\(\widetilde a,\widetilde c\in[0,1)\) are chosen. The exponential acts on those
representatives, not directly on the quotient classes.
Conversion to radians is introduced only upon choosing the
Archimedean chart
\[
 x:=A_{\rm rad}=\frac{\pi A}{180},\qquad
 y:=\Cstarrad=\frac{\pi\Cstar}{180},
\]
\[
 q_-=e^{-x+y},\qquad q_+=e^{-x-y}.
\]
The factor \(\pi/180\) belongs to this change analytic of coordinates. The
closure aritmetico of \(\alpha\), formulated before of this conversion, is
independent of that choice of chart angular.

\section{Central involutive algebra}

Consider the two-dimensional regular representation of the real algebra
generated by an involution. In a spectral basis, fix
\[
 R=\operatorname{diag}(1,-1),\qquad I=I_2,
\]
so that both eigenspaces have multiplicity one. Let
\[
 P_+=\frac{I+R}{2},\qquad P_-=\frac{I-R}{2}
\]
their spectral projectors. The central element determined by the
angular coordinates is
\[
 B_{A,\Cstar}=AI+\Cstar R
 =(A+\Cstar)P_++(A-\Cstar)P_-.
\]
The involution \(\Cstar\mapsto-\Cstar\) preserves the fixed projectors
\(P_\pm\) and interchanges the two spectral coefficients ---equivalently,
the associated channels---. Applying the exponential yields
\[
 T=e^{-xI-yR}=q_+P_++q_-P_-.
\]
Therefore, the two channels are the eigenvalues of a single element of the
algebra generated by \(I\) and \(R\).

\section{Invertibility of the channel chart}

\begin{theorem}[Reconstruction of the angular coordinates]
Para \(q_\pm>0\),
\[
 x=-\frac12\log(q_+q_-),
 \qquad
 y=\frac12\log\frac{q_-}{q_+}.
\]
Consequently, the map
\((A,\Cstar)\mapsto(q_+,q_-)\) is injective. The exchange
\(q_+\leftrightarrow q_-\) preserves \(A\) and transforms
\(\Cstar\mapsto-\Cstar\).
\end{theorem}
\begin{proof}
The multiplication and quotient of the two exponentials give
\(q_+q_-=e^{-2x}\) and \(q_-/q_+=e^{2y}\). Taking logarithms recovers
\(x\) and \(y\). Since \(\pi/180\neq0\), both angles are determined.
\end{proof}

\section{Addition and Recurrence Identities}

In this two-dimensional representation, the even and odd moments of the operator
\(T\) are
\[
 c_n=\Tr(T^n)=q_-^n+q_+^n,
 \qquad
 s_n=-\Tr(RT^n)=q_-^n-q_+^n.
\]

The multiplicity-one hypothesis is part of the domain of these formulas.
In a representation with multiplicities \(m_+\) and \(m_-\), the corresponding
traces would be
\(m_+q_+^n+m_-q_-^n\) and
\(m_-q_-^n-m_+q_+^n\), respectively.

\begin{theorem}[Laws of Composition of Moments]
For \(m,n\geq0\),
\[
 c_{m+n}=\frac{c_mc_n+s_ms_n}{2},
 \qquad
 s_{m+n}=\frac{s_mc_n+c_ms_n}{2},
\]
\[
 c_n^2-s_n^2=4e^{-2nx},
\]
and both sequences satisfy
\[
 X_{n+2}=c_1X_{n+1}-e^{-2x}X_n.
\]
Moreover,
\[
 \partial_Ac_n=-\frac{n\pi}{180}c_n,
 \quad
 \partial_As_n=-\frac{n\pi}{180}s_n,
\]
\[
 \partial_{\Cstar}c_n=\frac{n\pi}{180}s_n,
 \quad
 \partial_{\Cstar}s_n=\frac{n\pi}{180}c_n.
\]
\end{theorem}
\begin{proof}
Expand the products of
\(q_-^m\pm q_+^m\) with \(q_-^n\pm q_+^n\). The quadratic identity follows
from \((u+v)^2-(u-v)^2=4uv\). The recurrence corresponds to the characteristic
polynomial \(z^2-c_1z+q_-q_+\), and the derivatives are obtained using the
chain rule.
\end{proof}

\begin{corollary}[Determination of the complete hierarchy]
With \(A\) and \(\Cstar\) fixed, all pairs \((c_n,s_n)\) are
determined by the recurrence. No physical species introduces an additional
channel.
\end{corollary}

Finally,
\[
 s_n=2e^{-nx}\sinh(ny),
 \qquad
 c_n=2e^{-nx}\cosh(ny).
\]
The coordinate \(A\) determines the common decay, and \(\Cstar\) the oriented
separation. The decimal order of the action line is not introduced into these
formulas. It is deduced in \cref{chap:escala-accion} from the integral cokernel
of the local Hadamard block. The subsequent choice of a basis
expressed in \(\mathrm{J\,s}\) belongs to the metrological transducer and does not
alter the dimensionless identities.
